\specialsection{Заключение}

Выполненная работа показала, что схема отказоустойчивого хранения может
рассматриваться не как фиксированный выбор между репликацией и кодированием
стираний, а как автоматически управляемое состояние данных. Для данных с
дозаписью такой подход позволяет локально менять форму избыточности в фоне:
сохранять дешёвое чтение для более активных участков и компактнее хранить
менее активные данные. Поставленная во введении цель настоящей работы
достигнута: предложенный подход соответствует требованиям отказоустойчивости и
позволяет автоматически менять форму избыточности между состояниями с
преобладанием реплик и состояниями с преобладанием проверочных блоков в
зависимости от активности данных, занятого места, ожидаемой стоимости чтения и
стоимости фонового перехода между схемами.

В ходе работы были получены следующие результаты:
\begin{itemize}
    \item проанализированы существующие подходы к отказоустойчивому хранению:
    репликация, коды стираний, локально восстанавливаемые и конвертируемые
    коды, температурные политики и системы с изменением схемы хранения;
    \item сформулирована модель гибридного хранения и температуры данных как
    численной метрики текущей интенсивности чтений;
    \item определены допустимые переходы между схемами, включая обмен
    реплик на проверочные блоки, обратные переходы, изменение числа
    локальных групп, объединение и разделение совместимых схем;
    \item разработана модель оценки полезности перехода, учитывающая занятое
    место, ожидаемую стоимость чтения и стоимость фонового перекодирования;
    \item спроектирована архитектура адаптивного выбора переходов и разработана
    экспериментальная система \nolinkurl{https://github.com/v131v/hybrid_distributed_storage} с фоновым перекодировщиком, локальными
    планировщиками и глобальным планировщиком;
    \item проведена экспериментальная оценка поведения системы при охлаждении
    данных, неравномерной нагрузке чтения и отказах узлов.
\end{itemize}

Все поставленные задачи были выполнены. Экспериментальные данные показали, что
предложенный подход в рассмотренных сценариях позволяет одновременно более эффективно
использовать дисковое пространство по сравнению с исключительно репликацией и снижать
задержку операций чтения по сравнению с исключительно кодированием стираний.
Также показана способность адаптироваться к разным типам нагрузки на примере
случайных чтений и чтений, распределённых по закону Ципфа. При этом сохраняется
отказоустойчивость: в рассмотренных сценариях отказов запросы чтения
продолжали обрабатываться по пути чтения с восстановлением.

Ограничения работы связаны с выбранной постановкой и уровнем реализации.
Рассматривалась append-нагрузка: запись новых объектов и дозапись существующих
без изменения уже закрытых полос данных. Фоновое восстановление физически
утраченных блоков было отделено от задачи адаптивного перекодирования. Кроме
того, планировщик итеративно выбирает следующий лучший переход, что не
гарантирует глобального оптимума конфигурации. Предложенный подход также не
предназначен для мгновенной реакции на резкие изменения нагрузки.

Дальнейшее развитие работы может включать добавление механизма фонового
восстановления, проверку политики планирования на реальных распределениях
температуры данных, а также исследование масштабирования распределения
ответственности локальных планировщиков для крупных объектов и больших
кластеров. Тем не менее уже полученные результаты подтверждают основную
гипотезу: адаптивная смена схемы избыточности может быть полезным механизмом
управления способом хранения данных в отказоустойчивом хранилище.

\pagebreak
